\chapter{Theory}

\section{Information Processing Requirements}

Organizational information processing theory provides the frame through which
this thesis reads the task a retrieval system performs. The theory starts from
uncertainty, which \textcite[615]{tushman1978} treat as the difference between
the information required to complete a task and the information already
possessed. The tasks of organizational subunits vary in how uncertain they are,
and so in their information processing requirements. As work-related
uncertainty rises, so does the need for information, and with it the need for
capacity to process that information \parencite[615--616]{tushman1978}.
Requirements and capacity meet in the theory's central proposition: an
organization is more effective when the information processing capacity of its
structure matches the requirements it faces, and a mismatch in either direction
lowers performance. \textcite[619--620, 622]{tushman1978} present this
proposition as a hypothesis that remains to be fully tested, not as an
established result.

Structure supplies that capacity through coordination mechanisms.
\textcite[618]{tushman1978} order these mechanisms by their information
processing capacity, from rules and programs at the low end to formal
information systems and lateral relations at the high end. A technical system is
one means by which an organization raises its information processing capacity
\parencite[618]{tushman1978}. \textcite[29, 36]{galbraith1974} describes the
choice an organization faces as task uncertainty increases. The organization must either
reduce the amount of information it has to process or increase its capacity to
process it, and if it makes no conscious choice, reduced performance standards
follow automatically. Galbraith counts this lowering of performance standards
as the first strategy, and its cost shows up as overrun
schedules and budgets. For technology specifically, \textcite[65]{huber1990}
proposes that the use of advanced information technologies makes information
more available and more quickly retrieved. This proposition concerns
accessibility alone. Huber offers it, like his theory as a whole, as a
claim that still needs empirical investigation
\parencite[47, 67]{huber1990}.

In this thesis, the theory applies to the questions, not to the retriever.
Questions put to earnings call transcripts differ in their evidence
requirements: in what evidence an answer requires, and in where that evidence
sits. If effectiveness depends on a match between requirements and capacity, a
retrieval mechanism that treats every question as having the same evidence
requirement fails on some questions, however well it retrieves. The application
goes no further, because none of the sources above varies anything within a
technology: \textcite[67]{huber1990} states his propositions in general terms
and asks that hypotheses specify the particular technology, since even subtle
differences may count, and \textcite[622]{tushman1978} leave mechanisms other
than structure to future research. They provide no step from a change in chunk
size, retrieval strategy, or indexing representation to a change in information
processing capacity, so the hypotheses derived below, H1 to H3, rest on
retrieval mechanisms instead, because they concern properties of the retrieval
stage rather than of the organization. Where the theory does bear is on the
evaluation: whether a retrieved context suffices depends on what its question
requires of it, so a scheme that treats every question alike mismatches
requirement and capacity -- the reading the Discussion develops from the
partial-context result.

\section{Retrieval-Augmented Generation}

Retrieval-augmented generation pairs a retriever with a generator. The retriever
returns the passages it ranks highest for a question, and the generator produces
an answer conditioned on the question and those passages \parencite{lewis2020}.
An answer can therefore draw only on what was retrieved.
\textcite[1, 4]{kobeissi2026} formalize the generator as a function of the query
and the retrieved chunks alone. They argue that retrieval quality is the main
determinant of whether an answer is supported by the right context.

Two kinds of error follow from that structure: retrieval can fail to supply the
evidence an answer needs, or generation can misuse evidence it was given.
\textcite{kobeissi2026} argue that an answer-level score cannot separate the two,
since a wrong answer may come from either. Evaluation frameworks for these
systems accordingly score the retrieved context and the answer as distinct
properties \parencite{es2024,saad-falcon2024}.

On the generation side, the property at stake is faithfulness. Both frameworks
measure it as the grounding of an answer in the retrieved context. RAGAs asks
whether the claims an answer makes can be inferred from that context
\parencite{es2024}; ARES asks whether the answer adds nothing beyond it and
contradicts nothing in it \parencite{saad-falcon2024}. Defined this way,
faithfulness means grounding in the retrieved text, not factual truth, and this
thesis uses the term in that sense. Both frameworks assess it automatically, with
language-model judges \parencite{es2024,saad-falcon2024}.

The same definition fixes what a correct response is when the context cannot
support an answer. \textcite{muhamed2026} build their RefusalBench evaluation
around selective refusal: abstaining when defects in the provided context prevent
a reliable answer. RefusalBench scores an answer to such an instance as an error,
and a refusal over a minor, resolvable defect as a false refusal. When the
retrieved context cannot support an answer, abstaining is the correct behavior, and an
evaluation of answer quality has to credit it.

\section{Retrieval Design}

\subsection{Chunk Size}

The question of how much text should form a retrieval unit predates
retrieval-augmented generation. Research on passage retrieval treated the
definition and size of a passage as an open problem \parencite{callan1994}.
\textcite{kaszkiel1997} compared units ranging from paragraphs and sections to
fixed-length windows, and found no passage length that is always best.

In that work, unit size brings both an advantage and a cost. Smaller units dilute
evidence less: evidence concentrated in one region of a long document is not
averaged together with the unrelated text around it, which is why both studies
scored passages rather than whole documents \parencite{callan1994,kaszkiel1997}.
\textcite{callan1994} argues that small units also carry a cost. A block of
relevant text can be split across two of them, and a short passage matches fewer
of the terms in a long query. Both studies ranked documents by the passages they
contained; neither returned passages as retrieval units. Both measured length
in words, so they frame the question for this thesis without settling it.

Since the evidence supports both of the sizes compared here,
\thesisChunkSizeSmall{} and \thesisChunkSizeLarge{} tokens, the hypothesis
takes no direction. Retrieval quality is measured over the top five retrieved
chunks, by anchor coverage@5, the share of a question's gold anchors retrieved,
and MRR@5, the mean reciprocal rank of those anchors.

\begin{quote}
\textbf{H1:} Chunk size affects retrieval quality, measured as anchor coverage@5
and MRR@5.
\end{quote}

\subsection{Retrieval Strategy}

Lexical retrieval is represented here by BM25. As the best-known instance of the
probabilistic relevance framework, BM25 scores a document by weighting the query
terms it contains, so a term the document lacks contributes nothing
\parencite{robertson2009}. Dense retrieval encodes the question and each passage
independently and ranks passages by the dot product of their vectors
\parencite{karpukhin2020}.

Lexical and dense retrieval fail in different places. \textcite{karpukhin2020}
argue that they are complementary: term matching is sensitive to selective
keywords but poor at handling lexical variation, while dense retrieval captures
meaning but may miss rarely occurring key phrases. Reciprocal rank fusion
combines separate rankings into one. For each document, it sums the reciprocal
of a constant plus the document's rank in each ranking \parencite{cormack2009}.

The sources here do not settle whether any of these approaches is generally best.
Across a heterogeneous set of zero-shot retrieval tasks, BM25 remained a robust
baseline and no single approach consistently outperformed the others
\parencite{thakur2021}. That comparison included no fused system; the hybrid
strategy here fuses the dense and BM25 rankings by reciprocal rank fusion. The
hypothesis takes no direction.

\begin{quote}
\textbf{H2:} Retrieval strategy affects retrieval quality.
\end{quote}

\subsection{Indexing Representation}

A retriever can rank a passage only by the string it indexes. Metadata stored
with a passage but absent from that string, such as the company and period a
transcript belongs to, is invisible to lexical and dense retrieval alike,
however faithfully it is recorded.

In collections whose documents repeat one another, that absence matters.
\textcite[2]{yousuf2026} describe regulatory filings that reuse language across
companies and years: chunk-level semantic embeddings often cannot separate a
passage from its near-duplicates, and the metadata that would tell them apart is
usually applied only to filter results after retrieval. They instead serialize
structured fields into a header and concatenate it with the chunk text before
encoding \parencite[5]{yousuf2026}. They report that metadata-aware
representations of this kind improved on a plain-text baseline
\parencite[9]{yousuf2026}. In their field-level ablation, company and year
provide the strongest disambiguating signal, while removing section titles
yields only a modest drop at the chunk level and none at the document level
\parencite[12]{yousuf2026}. They present their work as, to their knowledge,
\enquote{the first systematic study showing that metadata signals significantly
improve retrieval on RAG systems} \parencite[3]{yousuf2026}.

The manipulation belongs to the line of document expansion, which enriches what
is indexed for a document instead of the query. In that line of work,
model-predicted queries are appended to the document \parencite{nogueira2019}.
The prefix here prepends recorded metadata from a fixed template, and its direct
precedent is the metadata-as-text construction above \parencite[5]{yousuf2026}.
The format differs: their header labels each field and is embedded by a dense
encoder only \parencite[5]{yousuf2026}; the prefix here carries no field
labels and is indexed by BM25 as well as by the dense encoder.

The hypothesis below therefore seeks corroboration, not a first test of an
untried mechanism, and extends their study where it stops. Their
chunking is fixed and their retrieval dense only \parencite[7--8]{yousuf2026},
while here indexing representation is crossed with both chunk size and retrieval
strategy. They report retrieval metrics alone and leave the evaluation of
generation quality to future work \parencite[8, 13]{yousuf2026}; this thesis
follows retrieval through to the answers. They also expect the benefit to carry
to other structured corpora, naming scientific articles, legal records, and
technical manuals \parencite[13]{yousuf2026}. Earnings call transcripts are
spoken and follow no template, so they lie outside that expectation. The same
banks discuss the same topics quarter after quarter, so the reasoning above
gives the hypothesis a direction.

\begin{quote}
\textbf{H3:} Enriching the indexed string with document metadata improves
retrieval quality.
\end{quote}

\section{From Retrieval to Answers}

Because an answer can draw only on what was retrieved, retrieval affects answers
through whether it supplies each question's evidence.
Where that evidence was retrieved, a faithful answer becomes possible. Where it
was not, abstention is the correct response. Evaluation frameworks for these
systems score the retrieved context and the answer for the same question
\parencite{es2024,saad-falcon2024}. \textcite{kobeissi2026} measure both stages
on one set of questions. None of these sources tests the link between the stages
question by question \parencite{es2024,saad-falcon2024,kobeissi2026}.

\textcite{yang2018} build a question-answering dataset around questions that
require information from more than one document to answer. The same retrieved
context can therefore be enough for one question and not for another: whether it
suffices depends on the question asked. Sufficiency is a property of a question
and a context together, not of the context alone.

{\centering\noindent Insert Figure \ref{fig:research-model} about here\par}

Figure~\ref{fig:research-model} sets out the research model: the three factors,
the hypothesis on each, and the answer stage examined beyond them.
